\documentclass[11pt, a4paper]{article}
\usepackage[a4paper, margin=2cm]{geometry}
\usepackage[utf8]{inputenc}
\usepackage[T2A]{fontenc}
\usepackage[english, russian]{babel}
\usepackage{amsmath, amssymb, amsthm}

\pagestyle{empty}

\begin{document}

\paragraph{№ 1}

\subparagraph{а)} 
Пусть $E_n = \left\{ x \in E \;\middle|\; |f(x)| \geqslant \frac{1}{n} \right\}$, $n \in \mathbb{N}$.

$$
0 = \int\limits_{E} |f(x)|\,d\mu(x) \geqslant \int\limits_{E_n} |f(x)|\,d\mu(x) \geqslant \int\limits_{E_n} \frac{1}{n}\,d\mu(x) = \frac{1}{n}\,\mu(E_n) \implies
$$

$$
\implies \frac{1}{n}\,\mu(E_n) \leqslant 0, \quad \mu(E_n) \geqslant 0, \quad \frac{1}{n} > 0 \implies \mu(E_n) = 0
$$

$$
\{ x \in E \mid |f(x)| > 0 \} = \bigcup_{n=1}^{\infty} E_n \leqslant \sum_{n=1}^{\infty} \mu(E_n) = 0 \implies f(x) = 0 \text{ п.в.}
$$

\subparagraph{б)}
$f = f^+ - f^-$

$A^+ = \{ x \in E \mid f(x) \geqslant 0 \}$, из (а): $f(\chi_{E \setminus A^+}) = 0$ п.в. $\implies \mu(\{x \in E \mid f(x) > 0\}) = 0$;

$A^- = \{ x \in E \mid f(x) < 0 \}$

$$
\int\limits_{A^-} f(x)\,d\mu = 0 \implies \int\limits_{A^-} (-f(x))\,d\mu = 0 \implies f(x) = 0 \text{ п.в. на } A^-
$$

\vspace{1.5em}
\hrule
\vspace{1.5em}

\paragraph{№ 2}

\subparagraph{1)}
$A = \{0, 1\}$, \quad $B = (0, 1)$

$$
d_x(A) = 0, \qquad d_{[x]}(B) = 0
$$
$$
d_x(B) = 1, \qquad d_{[x]}(A) = 1
$$

\subparagraph{2)}
$A = C$

$B = [0, 1] \setminus C$

$$
d_x(A) = 0 \qquad d_x(B) = 1
$$
$$
d_\lambda(B) = 0 \qquad d_\lambda(A) = 1
$$

\subparagraph{3)}
$A = \{0, 1\}$

$B = (0, 1)$

$$
d_\lambda(A) = d_\lambda(\{0\}) + d_\lambda(\{1\}) = 0
$$
$$
d_{[x]}(B) = 0
$$
$$
d_\lambda(B) = 1
$$
$$
d_{[x]}(A) = 1
$$

\vspace{1.5em}
\hrule
\vspace{1.5em}

\paragraph{№ 3}

$$
\mu_F((a, b]) = F(b) - F(a) = \int\limits_a^b F'(x)\,dx = \int\limits_{(a, b]} F'(x)\,dx
$$

\subparagraph{Теорема Леви:}
$(X, \mathcal{A}, \mu)$; $f_n \colon X \to [0, +\infty)$, $0 \leqslant f_1 \leqslant f_2 \leqslant \dots$ п.в. $x \in X$.

Если $f(x) = \lim\limits_{n \to +\infty} f_n(x)$ п.в. $x \in X$, то $f(x)$ --- измерима, и
$$
\int\limits_X f\,d\mu = \lim\limits_{n \to +\infty} \int\limits_X f_n\,d\mu
$$

\bigskip

\subparagraph{Вычисление интеграла:}

$$
\int\limits_{\mathbb{R}} f(x)\,d\mu_F(x) = \lim\limits_{n \to +\infty} \left( \int\limits_{\mathbb{R}} f_n(x)\,d\mu_F(x) \right) = \lim\limits_{n \to +\infty} \left( \int\limits_{\mathbb{R}} \sum_{k=1}^m c_k \mathbb{I}_{E_k}(x)\,d\mu(x) \right) =
$$

$$
= \lim\limits_{n \to +\infty} \left( \sum_{k=1}^m c_k \mu_F(E_k) \right) = \lim\limits_{n \to +\infty} \left( \sum_{k=1}^m c_k \int\limits_{E_k} F'(x)\,dx \right) =
$$

$$
= \lim\limits_{n \to +\infty} \left( \sum_{k=1}^m c_k \int\limits_{\mathbb{R}} \mathbb{I}_{E_k}(x) F'(x)\,dx \right) = \lim\limits_{n \to +\infty} \left( \int\limits_{\mathbb{R}} f_n(x) F'(x)\,dx \right) = \int\limits_{\mathbb{R}} f(x) F'(x)\,dx
$$

\vspace{1.5em}
\hrule
\vspace{1.5em}

\paragraph{№ 4}

$$
F_1(x) + C = F_2(x), \quad C \in \mathbb{R}
$$

\end{document}